\section{Discussion and Limitations}
\label{sec:discussion}

\todo{Write after v10 results. Planned points, kept only if the data
supports them:}
\begin{itemize}
  \item \textbf{Generator vs.\ curation.} Whether the pilot's ordering
    (easy generator $\gg$ curation $>$ none) survives calibration, i.e.\
    whether UED's benefit on 3D locomotion is mostly ``start easy''.
  \item \textbf{Regret estimates on locomotion.} Positive value loss
    depends on the critic, which here sees a privileged height scan;
    whether learnability is a better score for a continuous-control
    student with a binary success criterion.
  \item \textbf{Closed vs.\ open chain.} Whether the small closed-chain
    robot and Go1 favour the same methods and fail on the same features.
\end{itemize}

\paragraph{Limitations.}
\begin{itemize}
  \item Three seeds per method and robot; small differences (such as the
    pilot's tail gain of ACCEL-SFL) need more seeds to confirm.
  \item One student (brax PPO) and one budget per robot; methods that
    take fewer gradient steps (\plr, ACCEL) or simulate twice (PAIRED)
    are compared at equal environment steps, not equal compute.
  \item The level space is parametric and radially symmetric; it does not
    cover unstructured terrain maps or heading-dependent obstacles.
  \item The calibration uses pilot policies from the same benchmark, so
    the level space is tuned to this student and these robots.
  \item Simulation only: transfer is tested sim-to-sim (MJX to CPU
    MuJoCo), not on the physical Wojtek.
  \item PAIRED is run without the behaviour-cloning distillation
    of~\citet{mediratta2023stabilizing}, and ACCEL-SFL changes two
    switches relative to ACCEL.
\end{itemize}

\paragraph{Conclusion.}
\todo{Two sentences once the headline result is known.}
